% Circuit: GQSP for e^{-i delta A} via Motlagh-Wiebe Theorem 6.
% Matches notation in sec:prelim-QSP body text of 1_preliminaries.tex.
%
% Two dashed gategroups following the MW "last k" convention:
%   Box 1: d controlled-W calls (slots 1..d)
%   Box 2: d controlled-W^dag calls (slots d+1..2d)
% with a gap column between boxes to prevent visual overlap.
% PREP and PREP^dag shown as explicit gates on the b register.
%
\documentclass{article}
\usepackage[margin=1cm]{geometry}
\usepackage{tikz}
\usetikzlibrary{quantikz2}
\usepackage{amsmath,amssymb}
\usepackage{braket}

\newcommand{\ee}{\mathrm{e}}
\newcommand{\ii}{\mathrm{i}}

\tikzset{operator/.append style={rounded corners}}

\begin{document}
\pagestyle{empty}

\begin{figure}[ht]
  \centering
  \begin{quantikz}[row sep=0.7cm, column sep=0.5cm]
    \lstick{$\ket{0}_\mathrm{QSP}$}
      & \qw
      & \gate{R(\theta_0,\phi_0,\lambda_0)}
      & \octrl{1}
      \gategroup[3,steps=2,style={dashed,rounded corners,fill=black!6,inner sep=5pt},background,label style={label position=below,yshift=-0.45cm,align=center}]{\scriptsize Repeat $d$ times\\\scriptsize $j = 1,\dots,d$}
      & \gate{R(\theta_j,\phi_j,0)}
      & \qw
      & \octrl{1}
      \gategroup[3,steps=2,style={dashed,rounded corners,fill=black!6,inner sep=5pt},background,label style={label position=below,yshift=-0.45cm,align=center}]{\scriptsize Repeat $d$ times\\\scriptsize $j = d{+}1,\dots,2d$}
      & \gate{R(\theta_j,\phi_j,0)}
      & \qw
      & \rstick{$\bra{0}_\mathrm{QSP}$}\qw
      \\
    \lstick{$\ket{0}_b$}
      & \qwbundle{}
      & \gate{\textsc{Prep}}
      & \gate[2]{W}
      & \qw
      & \qw
      & \gate[2]{W^\dagger}
      & \qw
      & \gate{\textsc{Prep}^\dagger}
      & \rstick{$\bra{0}_b$}\qw
      \\
    \lstick{$\ket{\psi}$}
      & \qwbundle{}
      & \qw
      &
      & \qw
      & \qw
      &
      & \qw
      & \qw
      & \rstick{$\ee^{-\ii\delta A}\ket{\psi} + \mathcal{O}(\varepsilon_\mathrm{QSP})$}\qw
  \end{quantikz}
  \caption{Degree-$2d$ GQSP circuit realising $\ee^{-\ii\delta A}$
  from the walk operator $W$ of Eq.~\ref{eq:qsp-walk}.
  The first $d$ signal calls are controlled-$W$ and the last $d$
  are controlled-$W^\dagger$ (MW Thm.~6, $k = d$);
  post-selecting both ancillas on $\ket{0}$ yields
  $\ee^{-\ii\delta A} + \mathcal{O}(\varepsilon_\mathrm{QSP})$.}
  \label{circ:gqsp}
\end{figure}

\end{document}
